\chapter{ابنِ إجراءً يعتمد المنشآت الجديدة ويكلّف بالكشف}
\label{ch:walkthrough}

يضيف نور منشأة جديدة، فيبلّغ الإجراء أحمد ويطلب موافقته، ثم يكلّف نور بالكشف الأول أو يعيد الملف للاستكمال. تبنيه على أربع مراحل، وكل مرحلة تُنشر وتُجرَّب، ثم تتعلم الاختبار والتشغيل والإدارة حتى تتقن الإجراءات الداخلية.

\section{جهّز نفسك قبل البدء}
\label{sec:start}
\begin{rcquick}{جهّز نفسك قبل البدء}
\qphase{التحضير}
\qwho{أنت}{قبل الدخول إلى كيان}
\qwhere{هذه الخطوة}
\qdo{افتح «البيانات» واقرأ الكلمات الخمس؛ ستحتاجها من الخطوة التالية.}
\qdo{تعرّف إلى الفريق: أحمد يوافق، ونور يضيف المنشأة وينفّذ الكشف، وريم في تمرين الإتقان.}
\qdo{تدرّب بحساب واحد: اختر حسابك أنت في كل مكان يُذكر فيه أحمد أو نور.}
\qresult{تعرف ما سيحدث في نهاية المرحلة الأولى: يضيف نور منشأة، فيصل إلى أحمد إشعار.}
\qwhy{الإجراء الداخلي يعمل وحده خلف عمل الناس. لذلك تبنيه على مراحل صغيرة، وتجرّب كل مرحلة قبل التي بعدها، فتعرف أين وقع أي خطأ.}
\qdata{الإجراء الداخلي}{سلسلة خطوات يشغّلها النظام وحده عند حدث}{من «ادوات البناء» ← «الأساسية» ← «الإجراءات الداخلية».}
\qdata{المشغّل}{الحدث الذي يبدأ الإجراء}{مثل «إنشاء سجل جديد»؛ هو دائماً الخطوة 1.}
\qdata{الخطوة}{عمل واحد يقوم به الإجراء}{مثل «إرسال إشعار».}
\qdata{النشر}{يجعل ما بنيته يعمل فعلاً}{قبل النشر لا يبدأ الإجراء أبداً.}
\qdata{عملية التشغيل}{مرة واحدة عمل فيها الإجراء}{لكل منشأة جديدة عملية خاصة بها.}
\qdata{أحمد سالم — تجريبي}{رئيس فريق الكشف الميداني — 01}{يصله الإشعار ويرد على طلب الموافقة.}
\qdata{نور حسن — تجريبي}{باحث ميداني — 01}{يضيف المنشأة، وتصله المهمة أو طلب الاستكمال.}
\qdata{ريم خالد — تجريبي}{مدقق كشوف المنشآت — 01}{تُضاف في تمرين الإتقان.}
\qdata{الوحدة}{دائرة حماية المستهلك}{تتبع «مديرية حماية المستهلك وسلامة الغذاء» ثم «الإدارة العامة للتجارة الداخلية وحماية المستهلك».}
\end{rcquick}

\section{جهّز نوع السجلات الذي يبدأ الإجراء}
\label{sec:entity}
\begin{rcquick}{جهّز نوع السجلات الذي يبدأ الإجراء}
\qphase{التحضير}
\qwho{حساب المسؤول}{من يبني في «ادوات البناء»}
\qwhere{ادوات البناء ← الأساسية ← الكيانات ← إضافة كيان}
\qdo{«عنوان الكيان»: «المنشآت التجارية». «حقل 1» ← «اسم المنشأة» بنوع «نص».}
\qdo{«اضافة حقل جديد» بنوع «اختيار» باسم «حالة الملف»، و«إضافة اختيار» ثلاث مرات.}
\qdo{«انشاء كيان» ← في «إعدادات الكيان»: الحقل المعروض والتطبيق والوصف.}
\qdo{«انشاء كيان». لا تضف منشأة الآن.}
\qresult{«تم الإنشاء بنجاح»، و«المنشآت التجارية» في قائمة تطبيقك، وفي «إضافة» خانة الاسم وأزرار الحالات الثلاث.}
\qwhy{المشغّل يراقب نوع سجلات واحداً. حقل «حالة الملف» هو ما سيكتب فيه الإجراء قرار أحمد في المرحلة الرابعة.}
\qdata{عنوان الكيان}{المنشآت التجارية}{\req{} إن كان لديك من دليل الكشف، أضف «حالة الملف» فقط.}
\qdata{حقل 1}{اسم المنشأة}{نص.}
\qdata{حقل 2}{حالة الملف}{اختيار.}
\qdata{الاختيارات}{بانتظار الاعتماد، معتمدة للكشف، تحتاج استكمالاً}{\req{}}
\qdata{حقل الاسم المعروض}{اسم المنشأة}{}
\qdata{التطبيق الأساسي وقسمه}{تطبيق فريقك، مثل «حماية المستهلك»}{مع «إضافته إلى قائمة تطبيقه الأساسي».}
\qdata{الوصف}{ملفات المنشآت التي تتابعها الدائرة}{\req{} وإلا «صِف في 12 حرفاً وكلمتين على الأقل.»}
\end{rcquick}

\section{المرحلة 1: أنشئ الإجراء}
\label{sec:create}
\begin{rcquick}{المرحلة 1: أنشئ الإجراء}
\qphase{المرحلة 1: إشعار}
\qwho{حساب المسؤول}{من يبني الإجراء}
\qwhere{ادوات البناء ← الأساسية ← الإجراءات الداخلية ← إجراء جديد}
\qdo{في «الاسم» اكتب الاسم من «البيانات».}
\qdo{تحت «ما الذي يبدأه؟» اضغط بطاقة «إنشاء سجل جديد».}
\qdo{«إنشاء وفتح».}
\qresult{يفتح المحرر وفيه «إنشاء سجل جديد» ثم «إنهاء»، والرأس «لم يُنشر بعد» مع شارة صفراء «1 للإصلاح».}
\qwhy{نبدأ بأصغر إجراء ممكن: مشغّل وخطوة واحدة. إذا عمل، فكل ما نضيفه بعده يبني على شيء مضمون.}
\qdata{الاسم}{إشعار بمنشأة جديدة — تدريب}{\req{}}
\qdata{ما الذي يبدأه؟}{إنشاء سجل جديد}{\req{} مشغّلات أخرى في «مزيد من المُشغِّلات».}
\qshot{assets/screens/s1-new-sop.png}{تطلب اللوحة أمرين فقط: اسماً، وحدثاً يبدأ الإجراء. يبقى زر الإنشاء غير متاح حتى تختار الحدث. ① «الاسم»: اكتب اسماً يصف ما يفعله الإجراء. ② «إنشاء سجل جديد»: يبدأ الإجراء كلما أُضيف سجل في جدول تختاره بعد قليل. ③ «إنشاء وفتح».}
\end{rcquick}

\section{المرحلة 1: اختر الجدول الذي يراقبه المشغّل}
\label{sec:trigger}
\begin{rcquick}{المرحلة 1: اختر الجدول الذي يراقبه المشغّل}
\qphase{المرحلة 1: إشعار}
\qwho{حساب المسؤول}{من يبني الإجراء}
\qwhere{المحرر ← بناء ← الخطوة 1}
\qdo{اضغط بطاقة الخطوة 1 «عند: إنشاء سجل جديد».}
\qdo{خانة «الجدول»: اكتب «المنشآت» واختر «المنشآت التجارية».}
\qdo{لا يوجد زر حفظ؛ اقرأ «محفوظ» في الرأس.}
\qresult{تحت الخطوة 1 في «الخطوات» كلمة «جاهز»، والرأس «كل الخطوات جاهزة».}
\qwhy{«إنشاء سجل جديد» وحده لا يعرف أي سجلات تقصد، لذلك تحمل الخطوة «يحتاج إلى تهيئة» ولا يُنشر الإجراء قبل اختيار الجدول.}
\qdata{الجدول}{المنشآت التجارية}{\req{}}
\qshot{assets/screens/s1-trigger.png}{الصورة بعد اكتمال المراحل الأربع؛ ركّز الآن على الخطوة 1 فقط. ① «بناء»: الوضع الذي تعدّل فيه الخطوات. ② بطاقة المشغّل: اضغطها لتفتح إعدادها. ③ «الجدول»: اختر «المنشآت التجارية». ④ «الخطوات»: قائمة مرقمة تبيّن حالة كل خطوة، مثل «جاهز».}
\end{rcquick}

\section{المرحلة 1: أضف إشعاراً إلى أحمد}
\label{sec:notify}
\begin{rcquick}{المرحلة 1: أضف إشعاراً إلى أحمد}
\qphase{المرحلة 1: إشعار}
\qwho{حساب المسؤول}{من يبني الإجراء}
\qwhere{المحرر ← بناء ← «+» بين الخطوة 1 و«إنهاء»}
\qdo{اضغط «+»، واكتب في البحث «إرسال إشعار» واختره.}
\qdo{املأ «عنوان الرسالة» و«محتوى الرسالة» من «البيانات».}
\qdo{«المستلمون»: اكتب اسم أحمد (أو اسمك) واختره، ثم \en{Esc}.}
\qdo{اترك «الأولوية» «عادية» والبقية فارغة.}
\qresult{ثلاث بطاقات: «إنشاء سجل جديد» ← «إرسال إشعار» ← «إنهاء»، والرأس «كل الخطوات جاهزة».}
\qwhy{هذه أول خطوة يرى نتيجتها إنسان. نكتب الآن نصاً ثابتاً، وفي المرحلة الثانية نضيف إليه اسم المنشأة.}
\qdata{عنوان الرسالة}{أُضيفت منشأة جديدة}{\req{}}
\qdata{محتوى الرسالة}{افتح المنشآت التجارية لترى الملف الجديد وراجعه.}{}
\qdata{المستلمون}{أحمد، أو حسابك}{\req{} عضو محدد.}
\qdata{الأولوية}{عادية}{}
\qshot{assets/screens/s1-add-step.png}{في النافذة مئات الخطوات، والبحث بالاسم أسرع طريق إليها. خطوات كيان كلها في فئة «Kayan». ① خانة البحث: اكتب اسم الخطوة. ② «Kayan»: فئة خطوات النظام. ③ «إرسال إشعار»: اضغطها لتضيفها.}
\qshot{assets/screens/s1-notify.png}{الصورة بعد المرحلة الثانية، لذلك في العنوان اسم المنشأة مُدرجاً؛ الآن تكتب العنوان نصاً ثابتاً فقط. ① «عنوان الرسالة». ② «محتوى الرسالة». ③ «المستلمون»: العضو الذي يصله الإشعار.}
\end{rcquick}

\section{المرحلة 1: انشر الإصدار الأول}
\label{sec:publish}
\begin{rcquick}{المرحلة 1: انشر الإصدار الأول}
\qphase{المرحلة 1: إشعار}
\qwho{حساب المسؤول}{من يبني الإجراء}
\qwhere{المحرر ← مراجعة ونشر}
\qdo{اضغط «مراجعة ونشر».}
\qdo{في «الفحوصات» يجب أن ترى «كل الخطوات مهيّأة.». إن ظهرت مشكلة، اضغطها وأصلحها.}
\qdo{اضغط «نشر \en{v1}».}
\qresult{رسالة «أصبح \en{v1} منشورًا»، والرأس «منشور · \en{v1}» ثم «لا توجد تغييرات غير منشورة».}
\qwhy{ما تبنيه «مسودة» لا تعمل. النشر يصنع إصداراً مقفلاً يعمل؛ وأي تعديل بعده يصنع مسودة جديدة حتى تنشرها.}
\qdata{أول نشر}{\en{v1}}{\gen{} يزيد الرقم مع كل نشر.}
\qdata{بعد أي تعديل}{«المسودة v2 · تغييرات لم تُنشر بعد»}{لا يتغير ما يعمل حتى تنشر.}
\qshot{assets/screens/s1-publish.png}{تجمع اللوحة ما يجب أن تعرفه قبل النشر. في الصورة «نشر \en{v5}» لأنها بعد أربعة إصدارات؛ أول مرة «نشر \en{v1}». ① «مراجعة ونشر». ② «الفحوصات»: «كل الخطوات مهيّأة.» تعني أن لا شيء يمنع النشر. ③ زر النشر.}
\end{rcquick}

\section{المرحلة 1: أضف منشأة وشاهد الإجراء يعمل}
\label{sec:first-run}
\begin{rcquick}{المرحلة 1: أضف منشأة وشاهد الإجراء يعمل}
\qphase{المرحلة 1: إشعار}
\qwho{نور}{باحث ميداني — 01}
\qwhere{تطبيقك ← المنشآت التجارية ← إضافة}
\qdo{«اسم المنشأة»: من «البيانات».}
\qdo{«حالة الملف»: «بانتظار الاعتماد».}
\qdo{«إنشاء».}
\qresult{«فرن السنابل — تدريب» في القائمة بحالة «بانتظار الاعتماد».}
\qwhy{هذا هو الاختبار الحقيقي: يعمل نور كعادته، ولا يرى الإجراء، والإجراء يعمل خلفه لحظة الضغط على «إنشاء».}
\qdata{اسم المنشأة}{فرن السنابل — تدريب}{اسم جديد لكل تجربة.}
\qdata{التجارب التالية}{مخبز الريحان — تدريب، فرن القمح — تدريب، مطعم الياسمين — تدريب}{تميّز كل تشغيل وإشعار.}
\qshot{assets/screens/s1-add-record.png}{هذه شاشة الإضافة العادية، ولا شيء فيها يدل على الإجراء. ① «اسم المنشأة». ② «بانتظار الاعتماد»: حالة كل ملف جديد. ③ «إنشاء»: لحظة الضغط عليه يبدأ الإجراء.}
\end{rcquick}

\section{المرحلة 1: اقرأ عملية التشغيل والإشعار}
\label{sec:read-run}
\begin{rcquick}{المرحلة 1: اقرأ عملية التشغيل والإشعار}
\qphase{المرحلة 1: إشعار}
\qwho{حساب المسؤول، ثم أحمد}{من يبني الإجراء؛ المستلم}
\qwhere{المحرر ← عمليات التشغيل}
\qdo{اضغط «عمليات التشغيل»، ثم العملية «مكتمل».}
\qdo{على كل خطوة «نجح»، وتحت «التشغيل المحدد» عدد الخطوات الناجحة.}
\qdo{جرس «الإشعارات» ← تبويب «المستجدات».}
\qresult{إشعار «أُضيفت منشأة جديدة» في «المستجدات». لديك إجراء منشور يعمل.}
\qwhy{تتأكد من مكانين: عملية التشغيل تقول ماذا فعل كل خطوة، والإشعار يثبت أن النتيجة وصلت فعلاً.}
\qdata{حالات العملية}{في قائمة الانتظار، قيد التشغيل، بانتظار استجابة، مكتمل، فشل، ملغي}{}
\qdata{الإشعار غائب}{افتح «المستجدات» لا «للتنفيذ»}{إشعارات الإجراء هناك.}
\qshot{assets/screens/s1-runs.png}{عند اختيار عملية تظهر نتيجتها على اللوحة خطوة خطوة. الصورة من المرحلة الرابعة، فيها خطوات أكثر. ① «عمليات التشغيل». ② «التشغيل المحدد»: حالة العملية وعدد خطواتها الناجحة. ③ «نجح»: نتيجة الخطوة في هذه العملية.}
\qshot{assets/screens/s1-notification.png}{إشعارات الإجراء في «المستجدات»، لا في «للتنفيذ» المفتوح أولاً. ① جرس «الإشعارات». ② «المستجدات». ③ الإشعار: عنوانه ومحتواه كما كتبتهما.}
\end{rcquick}

\section{افهم المحرر: ثلاثة أوضاع ورأس يخبرك بالحالة}
\label{sec:editor}
\begin{rcquick}{افهم المحرر: ثلاثة أوضاع ورأس يخبرك بالحالة}
\qphase{افهم الأداة}
\qwho{حساب المسؤول}{من يبني الإجراء}
\qwhere{الإجراءات الداخلية ← إجراؤك}
\qdo{بدّل بين الأوضاع: «بناء» و«اختبار» و«عمليات التشغيل».}
\qdo{اقرأ الرأس: «محفوظ»، و«منشور · vN»، أو «المسودة vN · تغييرات لم تُنشر بعد».}
\qdo{افتح «سجل الإصدارات» وراجع إصداراتك المقفلة.}
\qresult{تعرف من الرأس وحده هل ما تراه منشور أم مسودة.}
\qwhy{الإصدارات المنشورة مقفلة، والتعديل يصنع مسودة. عمليات التشغيل الجارية تكمل على إصدارها الأصلي حتى لو نشرت إصداراً جديداً.}
\qdata{بناء}{تعدّل الخطوات}{الحفظ تلقائي.}
\qdata{اختبار}{«وضع الاختبار · محاكاة»}{تجرّب المسودة دون آثار.}
\qdata{عمليات التشغيل}{ما حدث فعلاً}{تصفية «فعلي» أو «اختبارات».}
\qdata{سجل الإصدارات}{«الإصدارات المنشورة مقفلة…»}{«عرض اللقطة» لرؤية إصدار قديم.}
\end{rcquick}

\section{المرحلة 2: أدرج اسم المنشأة في الإشعار}
\label{sec:insert-data}
\begin{rcquick}{المرحلة 2: أدرج اسم المنشأة في الإشعار}
\qphase{المرحلة 2: موافقة}
\qwho{حساب المسؤول}{من يبني الإجراء}
\qwhere{بناء ← إرسال إشعار ← عنوان الرسالة}
\qdo{ضع المؤشر بعد آخر كلمة واكتب «: ».}
\qdo{«إدراج بيانات» ← «ابحث في الحقول»: اكتب «فرن».}
\qdo{اضغط السطر \en{data.asm\_almnshat}.}
\qresult{العنوان «أُضيفت منشأة جديدة: \en{\{\{ \$event.data.asm\_almnshat \}\}}»، والرأس «المسودة \en{v2} · تغييرات لم تُنشر بعد».}
\qwhy{لا نكتب الاسم بأيدينا: الرمز بين القوسين المزدوجين يُستبدل باسم كل منشأة عند التشغيل. \en{\$event} هو ما حمله المشغّل.}
\qdata{\en{\$event.data}}{السجل الذي بدأ الإجراء}{حقوله بأسمائها الإنجليزية المولَّدة.}
\qdata{\en{\{\{ … \}\}}}{مرجع داخل نص}{يُستبدل بالقيمة عند التشغيل.}
\qdata{القائمة لا تعرض حقولك}{أضف منشأة مرة بعد النشر}{القائمة تعرض قيم آخر تشغيل أو اختبار.}
\qdata{تحتاج إلى صيغة؟}{زر الشيفرة \en{</>} بجانب الحقل}{لكتابة تعبير كامل.}
\qshot{assets/screens/s2-insert-data.png}{تعرض القائمة حقول آخر عملية تشغيل مع قيمها، فتعرف الحقل من قيمته لا من اسمه. ① «إدراج بيانات». ② «ابحث في الحقول»: اكتب اسم منشأة، أو \en{asm}. ③ \en{data.asm\_almnshat}: الاسم في السجل الجديد؛ اضغطه.}
\end{rcquick}

\section{المرحلة 2: أضف «طلب موافقة»}
\label{sec:approval}
\begin{rcquick}{المرحلة 2: أضف «طلب موافقة»}
\qphase{المرحلة 2: موافقة}
\qwho{حساب المسؤول}{من يبني الإجراء}
\qwhere{بناء ← «+» بين «إرسال إشعار» و«إنهاء»}
\qdo{ابحث عن «طلب موافقة» واختره.}
\qdo{«عنوان الموافقة»: اكتب النص، ثم «إدراج بيانات» ← اسم المنشأة.}
\qdo{«تفاصيل الموافقة»: من «البيانات».}
\qresult{الخطوة على اللوحة بين الإشعار و«إنهاء»، وعنوانها تحت اسمها.}
\qwhy{الإشعار يُعلم فقط. «طلب موافقة» يطلب قراراً ويوقف الإجراء حتى يرد أحمد، ثم يحفظ ما اختاره.}
\qdata{عنوان الموافقة}{اعتماد إدراج منشأة: + اسم المنشأة مُدرجاً}{\req{}}
\qdata{تفاصيل الموافقة}{راجع ملف المنشأة. وافق إذا كانت البيانات مكتملة، أو اطلب تغييرات واكتب الناقص.}{}
\qshot{assets/screens/s2-approval-fields.png}{العنوان ما يراه أحمد في قائمة طلباته، والتفاصيل تخبره بما يراجعه. ① «عنوان الموافقة»: مطلوب. ② «تفاصيل الموافقة».}
\end{rcquick}

\section{المرحلة 2: اختر من يوافق}
\label{sec:approver}
\begin{rcquick}{المرحلة 2: اختر من يوافق}
\qphase{المرحلة 2: موافقة}
\qwho{حساب المسؤول}{من يبني الإجراء}
\qwhere{طلب موافقة ← مراحل الموافقة}
\qdo{«إضافة مرحلة».}
\qdo{الخانة الأولى (فيها «مدير» رمادياً): اكتب \en{lead}.}
\qdo{«تحديد» ← اختر أحمد (أو نفسك) ← \en{Esc}.}
\qresult{اختفت الرسالتان «مفتاح المرحلة مطلوب» و«مطلوب معتمد واحد على الأقل».}
\qwhy{الطلب يمر بمراحل، وكل مرحلة شخص أو أكثر يجب أن يرد. «إضافة مرحلة» أخرى تضيف مستوى اعتماد بعد الأول.}
\qdata{مفتاح المرحلة}{\en{lead}}{\req{} إنجليزي قصير، لا يراه أحمد.}
\qdata{المعتمد}{أحمد، أو حسابك}{\req{} واحد على الأقل.}
\qshot{assets/screens/s2-approval-stage.png}{تحتاج كل مرحلة مفتاحاً ومعتمداً واحداً على الأقل. ① مفتاح المرحلة: \en{lead}. ② المعتمد. ③ «إضافة مرحلة».}
\end{rcquick}

\section{المرحلة 2: اكتب أزرار الرد}
\label{sec:actions}
\begin{rcquick}{المرحلة 2: اكتب أزرار الرد}
\qphase{المرحلة 2: موافقة}
\qwho{حساب المسؤول}{من يبني الإجراء}
\qwhere{طلب موافقة ← إجراءات الطلب}
\qdo{«الإجراء 1»: غيّر التسمية \en{Approve} إلى «موافقة»، وتحت «التوقيع» اضغط «لا شيء».}
\qdo{«الإجراء 2»: غيّر \en{Request changes} إلى «طلب تغييرات»، واترك «التعليق» «مطلوب».}
\qdo{لا تغيّر المفاتيح ولا النتائج.}
\qresult{«الإجراء 1» «موافقة» بلا توقيع، و«الإجراء 2» «طلب تغييرات».}
\qwhy{لكل زر تسمية يراها أحمد، ونتيجة يحفظها النظام ويقرؤها الإجراء بعد الرد. نغيّر التسمية فقط، لأن التفرّع سيقرأ النتيجة.}
\qdata{الإجراء 1}{مفتاح \en{approve}، تسمية «موافقة»، نتيجة \en{approved}}{التوقيع «لا شيء».}
\qdata{الإجراء 2}{المفتاح كما هو، تسمية «طلب تغييرات»، نتيجة \en{needs\_changes}}{التعليق «مطلوب».}
\qshot{assets/screens/s2-approval-actions.png}{لكل زر ثلاثة أشياء: ما يراه أحمد، وما يحفظه النظام، وما يُطلب منه عند الضغط. ① التسمية: «موافقة». ② النتيجة: \en{approved}. ③ «لا شيء» تحت «التوقيع»: الموافقة بلا توقيع.}
\end{rcquick}

\section{المرحلة 2: «مرفقات المكتبة» والمهلة}
\label{sec:attachments}
\begin{rcquick}{المرحلة 2: «مرفقات المكتبة» والمهلة}
\qphase{المرحلة 2: موافقة}
\qwho{حساب المسؤول}{من يبني الإجراء}
\qwhere{طلب موافقة ← مرفقات المكتبة}
\qdo{اضغط زر التعبير \en{</>} بجانب «مرفقات المكتبة».}
\qdo{احذف \en{""} واكتب \en{[} فقط؛ يكمل المحرر \en{[]}. ثم \en{Esc}.}
\qdo{«المهلة الزمنية»: غيّر \en{30 days} إلى \en{5 days}.}
\qresult{تحت «مرفقات المكتبة» «تعبير» و\en{[]}، و«المهلة الزمنية» \en{5 days}.}
\qwhy{في الإصدار الحالي يفشل الطلب إن بقيت «مرفقات المكتبة» فارغة. \en{[]} تعني «لا مرفقات». والمهلة هي مدة انتظار رد أحمد.}
\qdata{مرفقات المكتبة}{\en{[]}}{\req{} حل مؤقت لخلل في المنتج.}
\qdata{إن نسيتها}{«فشل عند 3 · طلب موافقة» ثم «Internal command failed with status 422»}{أضف \en{[]}، وانشر، وأضف منشأة جديدة.}
\qdata{المهلة الزمنية}{\en{5 days}}{بعدها ينتهي الطلب دون رد.}
\qshot{assets/screens/s2-attachments.png}{بعد الضغط على زر التعبير تتحول الخانة إلى تعبير، ويظهر بجانب عنوانها «تعبير». ① زر التعبير \en{</>}. ② \en{[]}.}
\end{rcquick}

\section{المرحلة 2: انشر، ثم دع أحمد يرد}
\label{sec:respond}
\begin{rcquick}{المرحلة 2: انشر، ثم دع أحمد يرد}
\qphase{المرحلة 2: موافقة}
\qwho{أحمد}{رئيس فريق الكشف الميداني — 01}
\qwhere{طلبات بانتظارك ← بانتظارك}
\qdo{المسؤول: «مراجعة ونشر» ← «نشر \en{v2}».}
\qdo{نور: أضف «مخبز الريحان — تدريب» بحالة «بانتظار الاعتماد».}
\qdo{أحمد: افتح «طلبات بانتظارك» ← بطاقة الطلب.}
\qdo{«موافقة» ← «تأكيد موافقة».}
\qresult{يخرج الطلب من «بانتظارك»، والعملية في «عمليات التشغيل» صارت «مكتمل» بعد أن كانت «بانتظار استجابة».}
\qwhy{الآن يتوقف الإجراء في منتصفه وينتظر إنساناً. أحمد لا يرى المحرر؛ يرى طلباً بعنوان واضح وزرين بالعربية.}
\qdata{رسالة بعد الإضافة}{«يتطلب «اعتماد إدراج منشأة: مخبز الريحان — تدريب» ردك.»}{تظهر أسفل الشاشة.}
\qdata{التعليق}{اختياري للموافقة}{إلزامي لطلب التغييرات.}
\qshot{assets/screens/s2-approval-request.png}{أحمد يرى طلباً بعنوان وتفاصيل، ومن يجب أن يرد، وزرين كما كتبناهما. ① «بانتظارك»: الطلبات التي تنتظر ردك. ② بطاقة الطلب. ③ «مسار الموافقة»: المرحلة الوحيدة، وفيها أحمد. ④ «موافقة». ⑤ «طلب تغييرات»: يطلب تعليقاً قبل التأكيد.}
\end{rcquick}

\section{المرحلة 3: أعطِ طلب الموافقة اسماً مرجعياً}
\label{sec:reference}
\begin{rcquick}{المرحلة 3: أعطِ طلب الموافقة اسماً مرجعياً}
\qphase{المرحلة 3: تفرّع}
\qwho{حساب المسؤول}{من يبني الإجراء}
\qwhere{طلب موافقة ← متقدم}
\qdo{اضغط «متقدم» ← «امنحها اسمًا مقروءًا».}
\qdo{في «الاسم المرجعي» احذف \en{step} واكتب \en{approval} ← \en{Enter}.}
\qresult{«الاسم المرجعي» \en{approval}.}
\qwhy{لتسأل خطوةٌ عن نتيجة خطوة قبلها تحتاج اسمها المرجعي: \en{\$steps.approval.output} هو ما أعادته خطوة \en{approval}.}
\qdata{الاسم المرجعي}{\en{approval}}{\req{} أحرف إنجليزية صغيرة وأرقام و\_ ، يبدأ بحرف.}
\qdata{\en{\$steps.<الاسم>.output}}{مخرجات خطوة سابقة}{تغيير الاسم يحدّث كل صيغة تشير إليه.}
\qshot{assets/screens/s3-reference.png}{الاسم المرجعي لا يظهر لأحد غيرك، لكنه طريقتك لتشير إلى نتيجة هذه الخطوة. ① «متقدم»: يفتح الإعدادات الإضافية. ② «الاسم المرجعي»: \en{approval}.}
\end{rcquick}

\section{المرحلة 3: أضف التفرّع وطريق «وافق»}
\label{sec:branch}
\begin{rcquick}{المرحلة 3: أضف التفرّع وطريق «وافق»}
\qphase{المرحلة 3: تفرّع}
\qwho{حساب المسؤول}{من يبني الإجراء}
\qwhere{بناء ← «+» بين «طلب موافقة» و«إنهاء»}
\qdo{ابحث عن «تفرّع شرطي» واختره. أبقِ «أول مسار مطابق».}
\qdo{في المسار \en{Branch 1}: «التسمية» «وافق».}
\qdo{زر التعبير بجانب «الشرط» ← احذف ما فيه والصق الشرط من «البيانات» ← \en{Esc}.}
\qresult{تحت المسار «شرط مخصص» والشرط، ورأس الخطوة «المسارات: وافق».}
\qwhy{الخطوة العادية يخرج منها طريق واحد. التفرّع يسأل سؤالاً لكل طريق، ويسلك أول طريق يصح شرطه.}
\qdata{المسارات المطابقة}{أول مسار مطابق}{}
\qdata{التسمية}{وافق}{تظهر على الخط في اللوحة.}
\qdoc{شرط «وافق»}{\en{\$steps.approval.output.terminalOutcomeKey == "approved"} \\ اقرأه: هل انتهى طلب خطوة \en{approval} بالنتيجة \en{approved}؟}
\qshot{assets/screens/s3-branch.png}{لكل مسار تسمية تظهر على الخط، وشرط يُسأل بعد رد أحمد. ① «أول مسار مطابق»: يسلك الإجراء أول طريق يصح شرطه. ② «التسمية»: «وافق». ③ الشرط بعد كتابته. ④ «إضافة مسار».}
\end{rcquick}

\section{المرحلة 3: أضف طريق «طلب تغييرات»}
\label{sec:branch-2}
\begin{rcquick}{المرحلة 3: أضف طريق «طلب تغييرات»}
\qphase{المرحلة 3: تفرّع}
\qwho{حساب المسؤول}{من يبني الإجراء}
\qwhere{تفرّع شرطي ← إضافة مسار}
\qdo{«إضافة مسار» ← «التسمية» «طلب تغييرات».}
\qdo{زر التعبير بجانب «الشرط» الثاني ← الصق الشرط ← \en{Esc}.}
\qdo{«ملاءمة سير العمل للعرض» لترى الإجراء كاملاً.}
\qresult{خطّان من «تفرّع شرطي»: «وافق» و«طلب تغييرات»، ولكل منهما «+» خاص.}
\qwhy{الشرط نفسه بنتيجة الزر الثاني. الطريقان يلتقيان عند «إنهاء».}
\qdoc{شرط «طلب تغييرات»}{\en{\$steps.approval.output.terminalOutcomeKey == "needs\_changes"}}
\qshot{assets/screens/s3-canvas.png}{الصورة بعد المرحلة الرابعة، فترى في كل طريق خطوتيه؛ الآن يكون كل طريق فارغاً. ① «تفرّع شرطي». ② طريق «وافق». ③ طريق «طلب تغييرات».}
\end{rcquick}

\section{المرحلة 4: مهمة لنور في طريق «وافق»}
\label{sec:task}
\begin{rcquick}{المرحلة 4: مهمة لنور في طريق «وافق»}
\qphase{المرحلة 4: التنفيذ}
\qwho{حساب المسؤول}{من يبني الإجراء}
\qwhere{بناء ← «+» على طريق «وافق»}
\qdo{ابحث عن «إنشاء مهمة» واختره.}
\qdo{«عنوان المهمة»: «كشف أول على: » ثم «إدراج بيانات» ← اسم المنشأة.}
\qdo{«المكلّفون» ← «تحديد» ← نور (أو نفسك) ← \en{Esc}.}
\qresult{على طريق «وافق» بطاقة «إنشاء مهمة»، وفي «الخطوات» «وافق · جاهز».}
\qwhy{الموافقة قرار، لكنها لا تنفّذ الكشف. المهمة تجعل الكشف عملاً له صاحب يراه في مهامه.}
\qdata{عنوان المهمة}{كشف أول على: + اسم المنشأة}{\req{} الحقل الإلزامي الوحيد.}
\qdata{المكلّفون}{نور، أو حسابك}{\rec{}}
\qshot{assets/screens/s4-task.png}{في «إنشاء مهمة» حقول كثيرة، والمطلوب منها العنوان فقط؛ نضيف المكلّف لتظهر عند نور. ① «عنوان المهمة». ② «المكلّفون».}
\end{rcquick}

\section{المرحلة 4: غيّر حالة الملف في طريق «وافق»}
\label{sec:update}
\begin{rcquick}{المرحلة 4: غيّر حالة الملف في طريق «وافق»}
\qphase{المرحلة 4: التنفيذ}
\qwho{حساب المسؤول}{من يبني الإجراء}
\qwhere{بناء ← «+» بعد «إنشاء مهمة»}
\qdo{ابحث عن «تحديث سجل كيان». «اسم الخطوة»: «اعتماد الملف».}
\qdo{«الكيان»: «المنشآت التجارية».}
\qdo{زر التعبير بجانب «السجل» ← \en{\$event.data.id} ← \en{Esc}.}
\qdo{«أضف قيمة الحقل» ← «حالة الملف» ← «معتمدة للكشف».}
\qresult{الخطوة «اعتماد الملف» «جاهز»، وفي «السجل» «المُشغِّل › data.id».}
\qwhy{من يفتح الملف يجب أن يرى أنه اعتُمد. \en{\$event.data.id} يعني «السجل الذي بدأ الإجراء»، فلا نعدّل منشأة ثابتة.}
\qdata{الكيان}{المنشآت التجارية}{\req{}}
\qdata{السجل}{\en{\$event.data.id}}{\req{} تعبير.}
\qdata{حالة الملف}{معتمدة للكشف}{}
\qshot{assets/screens/s4-update.png}{بعد كتابة التعبير يظهر في «السجل» «المُشغِّل › data.id»؛ الخطوة تعدّل المنشأة التي أُضيفت. ① «الكيان». ② «السجل»: المنشأة التي بدأت الإجراء. ③ القيمة الجديدة لـ«حالة الملف».}
\end{rcquick}

\section{المرحلة 4: جهّز طريق «طلب تغييرات»}
\label{sec:changes}
\begin{rcquick}{المرحلة 4: جهّز طريق «طلب تغييرات»}
\qphase{المرحلة 4: التنفيذ}
\qwho{حساب المسؤول}{من يبني الإجراء}
\qwhere{بناء ← «+» على طريق «طلب تغييرات»}
\qdo{أضف «إرسال إشعار» إلى نور بالقيم من «البيانات».}
\qdo{أضف بعده «تحديث سجل كيان» باسم «إعادة الملف للاستكمال».}
\qdo{«المنشآت التجارية»، و\en{\$event.data.id}، و«حالة الملف» = «تحتاج استكمالاً».}
\qresult{في «الخطوات» تسع خطوات «جاهز»، والرأس «كل الخطوات جاهزة».}
\qwhy{إذا طلب أحمد تغييرات يجب أن يعرف نور، وأن يرى الملف أنه ناقص. تكرر ما تعلمته بقيم مختلفة.}
\qdata{عنوان الرسالة}{استكمل ملف: + اسم المنشأة}{}
\qdata{محتوى الرسالة}{طلب رئيس الفريق تغييرات. افتح طلب الموافقة لتقرأ ملاحظته، ثم صحّح الملف.}{}
\qdata{المستلمون}{نور}{}
\qdata{حالة الملف}{تحتاج استكمالاً}{}
\end{rcquick}

\section{المرحلة 4: انشر وجرّب الطريقين}
\label{sec:final-run}
\begin{rcquick}{المرحلة 4: انشر وجرّب الطريقين}
\qphase{المرحلة 4: التنفيذ}
\qwho{أحمد ونور}{المعتمد؛ الباحث}
\qwhere{مراجعة ونشر ← طلبات بانتظارك}
\qdo{«مراجعة ونشر» ← «نشر».}
\qdo{أضف «فرن القمح — تدريب» و«مطعم الياسمين — تدريب».}
\qdo{«فرن القمح»: «موافقة» ← «تأكيد موافقة».}
\qdo{«مطعم الياسمين»: «طلب تغييرات» ← «العنوان التفصيلي ناقص.» ← أكّد.}
\qresult{«فرن القمح» «معتمدة للكشف» مع مهمة «كشف أول على: فرن القمح — تدريب»، و«مطعم الياسمين» «تحتاج استكمالاً» مع إشعار «استكمل ملف…». عمليتان «مكتمل».}
\qwhy{الطريق الذي لم يُجرَّب لا تعرف إن كان يعمل. منشأتان: واحدة لكل طريق.}
\qshot{assets/screens/s4-records.png}{كل منشأة تحمل حالة تحكي ما جرى لها؛ منشآت المراحل الأولى بقيت كما هي لأن الإجراء لم يكن يغيّر الحالة يومها. ① «معتمدة للكشف»: «فرن القمح — تدريب». ② «تحتاج استكمالاً»: «مطعم الياسمين — تدريب». ③ «بانتظار الاعتماد»: منشآت المراحل الأولى.}
\end{rcquick}

\section{اختبر دون أثر: «محاكاة (آمنة)»}
\label{sec:simulate}
\begin{rcquick}{اختبر دون أثر: «محاكاة (آمنة)»}
\qphase{الاختبار والتشغيل}
\qwho{حساب المسؤول}{من يبني الإجراء}
\qwhere{المحرر ← اختبار}
\qdo{افتح وضع «اختبار». في «بيانات الاختبار» الصق الكائن من «البيانات».}
\qdo{اختر «محاكاة (آمنة)» ← «تشغيل المحاكاة».}
\qdo{اقرأ نتيجة كل خطوة على اللوحة. غيّر القيم وأعد المحاولة.}
\qresult{ترى أي طريق سلكه الإجراء دون إنشاء سجل أو إرسال شيء.}
\qwhy{المحاكاة تشغّل المسودة دون آثار، والخطوات التي تتصل بأنظمة أخرى تعيد بيانات نموذجية؛ وخطوة «تشغيل JavaScript» لا تعطي نتيجتها الحقيقية فيها.}
\qdata{محاكاة (آمنة)}{تشغّل المسودة}{لا سجلات ولا رسائل.}
\qdata{تشغيل فعلي}{يشغّل الإصدار المنشور}{«تشغيل vN فعليًا» بعد تأكيد الموافقة؛ له آثار حقيقية.}
\qdata{«انشر الإجراء أولًا…»}{لا يوجد إصدار منشور}{التشغيل الفعلي لا يعمل على المسودة.}
\qdata{بيانات غير صالحة}{«أصلح بيانات الاختبار: يجب أن تكون كائن JSON.»}{}
\qdoc{بيانات الاختبار}{\en{\{"data": \{"id": "test-1", "asm\_almnshat": "مخبز التجربة — تدريب"\}\}}}
\end{rcquick}

\section{اختبر خطوة واحدة}
\label{sec:test-step}
\begin{rcquick}{اختبر خطوة واحدة}
\qphase{الاختبار والتشغيل}
\qwho{حساب المسؤول}{من يبني الإجراء}
\qwhere{بناء ← تفرّع شرطي ← اختبار هذه الخطوة}
\qdo{اضغط خطوة «تفرّع شرطي».}
\qdo{«اختبار هذه الخطوة» بقيمة \en{approved} لمخرج \en{approval}، ثم بقيمة \en{needs\_changes}.}
\qdo{قارن المسار الذي اختاره مع ما توقعته.}
\qresult{\en{approved} تسلك «وافق»، و\en{needs\_changes} تسلك «طلب تغييرات».}
\qwhy{اختبار الخطوة وحدها أسرع طريقة لفحص شرط، دون انتظار رد معتمد. خطوات السجلات قد تطلب موافقتك لأنها تنفّذ فعلاً.}
\qdata{اختبار هذه الخطوة فقط}{يستخدم بيانات اختبار محفوظة}{}
\qdata{تحذير التنفيذ الفعلي}{«أُدرك أن اختبار هذه الخطوة سينفّذ إجراءً فعلياً…»}{يظهر في خطوات تغيّر البيانات.}
\end{rcquick}

\section{اقرأ تشغيلاً فاشلاً وأصلحه}
\label{sec:failed-run}
\begin{rcquick}{اقرأ تشغيلاً فاشلاً وأصلحه}
\qphase{الاختبار والتشغيل}
\qwho{حساب المسؤول}{من يبني الإجراء}
\qwhere{المحرر ← عمليات التشغيل ← فعلي}
\qdo{اختر العملية «فشل». اقرأ تحت «التشغيل المحدد» رقم الخطوة واسمها ورسالة الخطأ.}
\qdo{اضغط الخطوة الفاشلة واقرأ ما استلمته وما أعادته.}
\qdo{أصلح، ثم انشر، ثم شغّل حالة جديدة.}
\qdo{للخطوات المهمة: «متقدم» ← «إذا فشلت هذه الخطوة» واختر السلوك.}
\qresult{تعرف من العملية وحدها أين توقفت ولماذا.}
\qwhy{كل عملية تحفظ ما دخل كل خطوة وما خرج منها. الإصلاح لا يغيّر العمليات القديمة؛ ينطبق على ما يبدأ بعد النشر.}
\qdata{إيقاف التشغيل}{الافتراضي}{تُعلَّم العملية فاشلة وتتوقف.}
\qdata{اتباع مسار الفشل}{مساران: نجاح وفشل}{تستمر الأخطاء في مسار الفشل.}
\qdata{المتابعة}{يكمل إلى الخطوة التالية}{تقرأ الخطوات اللاحقة الخطأ.}
\qdata{إعادة المحاولة قبل الاستسلام}{يعيد المحاولة تلقائياً}{للأخطاء العابرة.}
\qdata{لا عملية تشغيل أصلاً}{غير منشور، أو جدول آخر في المشغّل}{}
\qdata{عدّلت ولم يتغير شيء}{لم تنشر بعد التعديل}{}
\qdata{سلك طريقاً غير متوقع}{اقرأ \en{terminalOutcomeKey} في خطوة الطلب}{قارنه حرفاً بحرف بالشرط.}
\qdata{تنتظر ولا تكتمل}{«بانتظار استجابة»}{طبيعي حتى يرد المعتمد أو تنتهي المهلة.}
\end{rcquick}

\section{أدِر إجراءاتك}
\label{sec:manage}
\begin{rcquick}{أدِر إجراءاتك}
\qphase{الاختبار والتشغيل}
\qwho{حساب المسؤول}{من يبني الإجراء}
\qwhere{ادوات البناء ← الأساسية ← الإجراءات الداخلية}
\qdo{استخدم التصفية: «منشورة» و«مسودات» و«تحتاج إلى متابعة» و«متوقفة مؤقتًا» و«مؤرشفة».}
\qdo{لإجراء يفشل: «فتح التشغيل الفاشل».}
\qdo{لإيقافه مؤقتاً: «إيقاف الإجراء مؤقتًا»، ثم «استئناف».}
\qresult{ترى من القائمة صحة كل إجراء: «عمليات التشغيل الأخيرة الناجحة» وآخر فشل.}
\qwhy{الإجراء المنشور يعمل وحده أسابيع، فيحتاج نظرة دورية. الإيقاف المؤقت يمنع بدء تشغيلات جديدة دون حذف شيء.}
\qdata{إيقاف الإجراء مؤقتًا}{لا يبدأ أي تشغيل جديد}{حتى «استئناف».}
\qdata{تحتاج إلى متابعة}{إجراءات فيها فشل حديث}{}
\qdata{المجلدات}{رتّب الإجراءات بحسب الفريق}{مثل «حماية المستهلك — تدريب».}
\end{rcquick}

\section{تعرّف إلى بقية المشغّلات والخطوات}
\label{sec:catalog}
\begin{rcquick}{تعرّف إلى بقية المشغّلات والخطوات}
\qphase{افهم الأداة}
\qwho{حساب المسؤول}{من يبني الإجراء}
\qwhere{إجراء جديد ← مزيد من المُشغِّلات، ثم «+» في المحرر}
\qdo{افتح «مزيد من المُشغِّلات» واقرأ وصف كل مشغّل في «البيانات». لا تنشئ شيئاً.}
\qdo{في إجراء تدريبي اضغط «+» وابحث عن كل خطوة في الجدول.}
\qdo{اربط كل واحد بمثال من عملك.}
\qresult{تعرف أي مشغّل تختار لكل حدث، وأي خطوة لكل عمل.}
\qwhy{الإجراء الذي بنيته استخدم مشغّلاً وست خطوات فقط. هذه بقية الأدوات التي تحتاجها لأغلب الإجراءات.}
\qdata{تحديث سجل}{عند تحديث سجل في جدول}{يحمل \en{\$event.data} و\en{\$event.oldData} و\en{\$event.changedFields}؛ المثال الثاني.}
\qdata{حذف سجل}{عند حذف سجل}{}
\qdata{جدولة}{وفق جدول زمني}{تقرير أسبوعي مثلاً.}
\qdata{يدوي}{يشغّله عضو بنفسه}{مدخلاته في \en{\$event.input}.}
\qdata{عند استدعاء النموذج}{يستدعيه نموذج ويرد عليه}{ينتهي بـ«إرجاع استجابة النموذج».}
\qdata{عند إرسال نموذج}{عند إرسال جلسة نموذج}{}
\qdata{استدعاء من إجراء آخر}{يستقبل مدخلات «استدعاء إجراء تشغيلي» في \en{\$event}}{}
\qdata{جلب سجل}{يقرأ سجلاً كما هو الآن}{المثال الثاني.}
\qdata{البحث عن سجلات كيان}{يجد عدة سجلات بشروط}{}
\qdata{تأخير}{ينتظر مدة ثم يتابع}{المثال الثاني.}
\qdata{لكل عنصر}{يكرر مساراً لكل عنصر من قائمة}{داخله \en{\$item} و\en{\$index}.}
\qdata{تشغيل JavaScript}{شيفرة في بيئة معزولة}{بلا شبكة؛ متغيراتها بـ\en{context.variables.NAME}.}
\qdata{طلب نموذج}{يطلب من الأعضاء إكمال نموذج}{}
\qdata{إرسال تذكير بالموافقة}{يذكّر المعتمد بطلب معلّق}{}
\qdata{المتغيرات}{قيم ثابتة يستخدمها الإجراء في التعبيرات}{القيم المشفّرة لا تظهر مرة أخرى.}
\qdata{الصلاحيات}{عضو التشغيل يحتاج قراءة السجل}{والمستلمون والمعتمدون أعضاء بتعيين نشط.}
\end{rcquick}

\section{طبّق وحدك: إشعار ريم بعد الموافقة}
\label{sec:practice}
\begin{rcquick}{طبّق وحدك: إشعار ريم بعد الموافقة}
\qphase{الإتقان}
\qwho{أنت}{دون الرجوع إلى الخطوات}
\qwhere{الإجراء نفسه ← بناء}
\qdo{قرر: في أي طريق تضيف الإشعار؟ ولماذا ليس قبل التفرّع؟}
\qdo{أضف «إرسال إشعار» إلى ريم مع اسم المنشأة مُدرجاً، ثم انشر.}
\qdo{أضف منشأة ووافق عليها: خطوة ريم «نجح»، وخطوات «طلب تغييرات» لم تعمل.}
\qdo{أضف منشأة أخرى واطلب تغييرات: لم يصل ريم شيء.}
\qresult{وضعت الخطوة في المكان الصحيح من أول مرة، وفسّرت النتيجتين من «عمليات التشغيل».}
\qwhy{تتقن الإجراءات عندما تبني واحداً دون أن تنظر إلى الخطوات، وتتنبأ بالطريق قبل الضغط.}
\qdata{المكان}{طريق «وافق»}{قبل التفرّع يصل الإشعار في الحالتين.}
\qdata{المستلم}{ريم خالد — تجريبي}{أو حسابك.}
\end{rcquick}

\section{اختبر فهمك}
\label{sec:quiz}
\begin{rcquick}{اختبر فهمك}
\qphase{الإتقان}
\qwho{أنت}{أجب قبل أن تفتح «البيانات»}
\qwhere{الإجراء نفسه}
\qdo{أجب عن كل سؤال بصوت عالٍ.}
\qdo{افتح «البيانات» وقارن جوابك.}
\qdo{تحقق من كل جواب في المكان المذكور في كيان.}
\qresult{أجبت عن الأسئلة كلها، وتعرف أين ترى كل جواب في كيان.}
\qwhy{الإتقان أن تشرح لماذا حدث ما حدث، لا أن تتذكر مواضع الأزرار.}
\qdata{عدّلت نص الإشعار ولم يتغير ما يصل؛ لماذا؟}{لم تنشر بعد التعديل}{الرأس: «تغييرات لم تُنشر بعد».}
\qdata{ما الفرق بين «محاكاة (آمنة)» و«تشغيل فعلي»؟}{المحاكاة للمسودة بلا آثار؛ الفعلي للمنشور بآثار}{وضع «اختبار».}
\qdata{ماذا يعني \en{\$event.data.id}؟}{معرّف السجل الذي بدأ الإجراء}{«السجل» في «تحديث سجل كيان».}
\qdata{كيف تقرأ خطوة نتيجة خطوة قبلها؟}{\en{\$steps.<الاسم المرجعي>.output}}{«متقدم» ← «الاسم المرجعي».}
\qdata{عملية «بانتظار استجابة» منذ يومين: معطّلة؟}{لا؛ تنتظر المعتمد أو «المهلة الزمنية»}{«عمليات التشغيل».}
\qdata{لماذا لا نكتب اسم المنشأة بأيدينا؟}{«إدراج بيانات» يضع اسم كل منشأة في طلبها}{العنوان في «إرسال إشعار».}
\qdata{نشرت إصداراً جديداً؛ ماذا عن العمليات الجارية؟}{تكمل على إصدارها الأصلي}{«سجل الإصدارات».}
\qdata{أي طريق يسلك التفرّع إذا صح شرطان؟}{الأول، مع «أول مسار مطابق»}{إعداد «المسارات المطابقة».}
\qdata{خطوة تفشل أحياناً لخطأ عابر؛ ماذا تفعل؟}{«إعادة المحاولة قبل الاستسلام» أو «اتباع مسار الفشل»}{«إذا فشلت هذه الخطوة».}
\qdata{كيف توقف إجراءً دون حذفه؟}{«إيقاف الإجراء مؤقتًا»}{قائمة الإجراءات.}
\qdata{أين تجد إشعار الإجراء؟}{الجرس ← «المستجدات»}{لا «للتنفيذ».}
\end{rcquick}
